\section{INTRODUCCIÓN}

\subsection{Antecedentes}

La vigilancia de patógenos virales en aguas residuales se ha consolidado como una estrategia complementaria para observar la circulación de agentes infecciosos a escala poblacional pues, a diferencia de la vigilancia clínica, no depende del acceso a servicios de salud y de la notificación de casos. El análisis de aguas servidas integra señales biológicas de una comunidad conectada a una red de alcantarillado, enfoque que se denomina epidemiología basada en aguas residuales. Esto permite generar información agregada sobre presencia, persistencia y cambios temporales de microorganismos de interés sanitario \cite{Sims_2020}.

Durante la pandemia de COVID-19, la detección de ARN de SARS-CoV-2 en aguas residuales demostró que este tipo de muestra puede reflejar la circulación comunitaria del virus, incluso cuando la vigilancia clínica tiene limitaciones de cobertura o retrasos en la notificación \cite{Ahmed_2020,Medema_2020}. Además, la incorporación de métodos de secuenciación permitió pasar de la detección molecular hacia la vigilancia genómica, con capacidad para identificar linajes, variantes y diversidad viral presentes en una población \cite{Crits_Christoph_2021}. Este cambio amplió el valor de las aguas residuales como fuente de información para salud pública, especialmente en contextos urbanos donde múltiples sectores de la población convergen en una misma planta de tratamiento.

Además de SARS-CoV-2, las aguas residuales han sido utilizadas para la vigilancia ambiental de otros virus de interés sanitario, tales como poliovirus. Esta herramienta ha sido empleada dentro de los esfuerzos globales de erradicación para complementar la vigilancia de parálisis flácida aguda y detectar circulación silenciosa en comunidades \cite{Asghar_2014}. De forma similar, estudios recientes han mostrado que el ARN del virus sincitial respiratorio puede detectarse en matrices de aguas residuales y relacionarse con tendencias clínicas de infección respiratoria \cite{Hughes_2022}. Estos antecedentes sostienen la importancia de evaluar simultáneamente virus que puedan ser detectados en aguas residuales por la excreción de material viral por individuos infectados.

\subsection{Problemática}

La vigilancia de virus respiratorios y entéricos basada en casos clínicos depende de la presencia de síntomas, la búsqueda de atención, el acceso a diagnóstico y la notificación oportuna. Las revisiones sobre vigilancia en aguas residuales superan limitaciones recurrentes de ese circuito: cobertura poblacional incompleta, disponibilidad limitada de pruebas, retrasos de reporte y subregistro de infecciones leves o asintomáticas \cite{Bubba_2023,Parkins_2023}. En COVID-19, los casos notificados subestimaron la carga real por la restricción de pruebas y por la reducción del rol del testeo clínico durante la ola Ómicron BA.1, cuando la transmisión superó la capacidad diagnóstica disponible \cite{Parkins_2023}. Esta limitación disminuye la oportunidad de detectar circulación comunitaria y reduce la priorización de respuestas de salud pública.

La vigilancia basada en aguas residuales aporta una señal importante del área conectada al alcantarillado o planta de tratamiento. En Bangkok, en una investigación publicada en 2022, se analizaron 132 muestras de 19 plantas de tratamiento y detectaron ARN de SARS-CoV-2 cuando los casos reportados permanecían bajos; las cargas virales y los porcentajes de positividad se correlacionaron con los casos nuevos con rezagos de 22 a 24 días y valores de Spearman entre 0,85 y 1,00; lo cual indica una correlación positiva muy fuerte \cite{Sangsanont_2022}. En Calgary, investigadores estudiaron tres plantas de tratamiento y seis vecindarios durante un año, con detección de SARS-CoV-2 en 406 de 414 muestras (98,06 \%) y correlación con casos clínicos a escalas urbana, por planta y barrial \cite{Acosta_2022}. En países de ingresos bajos y medios, estos programas enfrentan limitaciones de infraestructura sanitaria, financiamiento, frecuencia de muestreo y capacidad para vigilancia de variantes \cite{Parkins_2023}.

La vigilancia clínica de poliovirus se apoya en la detección de parálisis flácida aguda, un desenlace que afecta a una fracción mínima de las personas infectadas: la tasa de ataque clínico del poliovirus virulento puede ser de un caso paralítico por cada 100 infecciones o menos \cite{Zurbriggen_2008}. En África, la confirmación de VDPV2 en 525 muestras de heces tomó una mediana de 49 días desde el inicio de la parálisis, y los brotes nuevos se confirmaron a los 35 días, tiempos vinculados con recolección, transporte, cultivo celular, RT-qPCR intratípica y secuenciación \cite{Shaw_2022}. En 2024, la detección de cVDPV2 en aguas residuales de Finlandia, Alemania, Polonia, España y Reino Unido mostró cepas genéticamente relacionadas y respaldó la vigilancia ambiental como mecanismo de detección temprana de reintroducción o circulación \cite{Bottcher_2025}.

El virus sincitial respiratorio (RSV) plantea otro problema de vigilancia por su carga en lactantes, adultos mayores e inmunocomprometidos, y por su presentación clínica compartida con otros virus respiratorios. Hughes y colaboradores detectaron ARN de RSV en sólidos sedimentados de aguas residuales de dos plantas y reportaron asociación fuerte con tasas de positividad clínica (Kendall tau = 0,65--0,77; $p < 10^{-7}$) \cite{Hughes_2022}. Le et al. desarrollaron un flujo de secuenciación ONT para muestras clínicas y aguas residuales, obtuvieron genomas completos en el 98 \% de las muestras con más de 10.000 copias/mL y observaron que las frecuencias de clados en aguas residuales representaban las frecuencias observadas en pacientes \cite{Le_2025}. El mismo estudio subraya la necesidad de vigilancia genómica accesible para contextos con recursos limitados, debido a la menor representación de datos genómicos de RSV procedentes de países de ingresos bajos y medios \cite{Le_2025}.

En Quito, la PTAR Quitumbe ofrece un punto concreto para observar esta brecha. Mejía Calle realizó un muestreo pasivo semanal de 24 horas en el influente de la PTAR Quitumbe, documentó un caudal de 75 L/s, cobertura de once barrios del Distrito Metropolitano de Quito y 47 muestras analizadas \cite{MejiaCalle_2024}. En ese trabajo, los genes N1 y N2 de SARS-CoV-2 fueron detectados en todas las muestras, y la secuenciación con Oxford Nanopore permitió identificar linajes de Ómicron, incluida la detección de JN.1 en aguas residuales antes de su identificación clínica local \cite{MejiaCalle_2024}. La evidencia local confirma la factibilidad del sitio para vigilancia ambiental, pero permanece concentrada en SARS-CoV-2. La brecha específica que se detecta en este TFM es la falta de información molecular y genómica multipatógeno en aguas servidas de Quito que integre SARS-CoV-2, poliovirus y RSV en una misma matriz ambiental de relevancia sanitaria \cite{MejiaCalle_2024,Asghar_2014,Hughes_2022}.

\subsection{Justificación}

La vigilancia en aguas residuales es una herramienta flexible y de bajo costo relativo para monitorear brotes y reemergencia de patógenos, con mayor sostenibilidad cuando integra varios blancos virales en un mismo sistema \cite{Bubba_2023}. En la PTAR Quitumbe, una muestra ambiental representa la circulación viral de los barrios conectados al alcantarillado. Estudiar SARS-CoV-2, poliovirus y RSV en una misma matriz reúne tres dimensiones sanitarias: el seguimiento de variantes respiratorias, la detección de circulación silenciosa de poliovirus y la observación de un virus respiratorio asociado con enfermedad grave en lactantes, adultos mayores e inmunocomprometidos \cite{Asghar_2014,Hughes_2022,Le_2025}.

En contextos con recursos limitados, los programas de vigilancia en aguas residuales enfrentan restricciones de infraestructura, financiamiento, frecuencia de muestreo y capacidad para monitorear variantes, aunque los sitios centinela de alcantarillado aportan información local útil para la toma de decisiones \cite{Parkins_2023}. Dos limitaciones documentadas refuerzan el valor de ampliar esta vigilancia: la confirmación de brotes de VDPV2 se retrasa cuando depende de muestras clínicas y circuitos de laboratorio complejos, y la vigilancia genómica de RSV mantiene baja representación de datos de países de ingresos bajos y medios \cite{Shaw_2022,Le_2025}. La PTAR Quitumbe ya demostró utilidad para SARS-CoV-2 y ofrece una base para avanzar hacia un enfoque multipatógeno \cite{MejiaCalle_2024}.

El estudio parte de un sitio ya caracterizado y de métodos aplicables a matrices ambientales complejas. La PTAR Quitumbe recibe aguas residuales de once barrios del Distrito Metropolitano de Quito, y un estudio previo aplicó muestreo pasivo semanal de 24 horas, obtuvo 47 muestras y secuenció 36 de ellas mediante Oxford Nanopore \cite{MejiaCalle_2024}. Existen además métodos dirigidos y metagenómicos enriquecidos para recuperar información genómica de virus humanos en aguas residuales, junto con protocolos de secuenciación aplicables a RSV y poliovirus \cite{Child_2023,Le_2025,Asghar_2014}. Sobre esa base, el trabajo abarca la detección molecular, la exploración genómica y la evaluación de las limitaciones propias de una matriz ambiental urbana.

Los beneficiarios directos son los equipos académicos y técnicos de microbiología ambiental, biología molecular, bioinformática y gestión de aguas residuales, que dispondrán de evidencia local sobre blancos virales, rendimiento de detección y limitaciones analíticas en la PTAR Quitumbe. Las instituciones de salud pública y saneamiento pueden usar estos insumos para valorar sitios centinela, priorizar patógenos y planificar vigilancia ambiental en Quito \cite{Parkins_2023,MejiaCalle_2024}. De forma indirecta, la población conectada a la PTAR y otros sectores urbanos se beneficiarían de futuros sistemas de monitoreo, en particular los grupos vulnerables ante RSV y las comunidades expuestas a circulación silenciosa de poliovirus o a cambios de variantes de SARS-CoV-2 \cite{Hughes_2022,Bottcher_2025,Crits_Christoph_2021}.

\subsection{Problema de investigación, pregunta y objetivos}

El vacío local que aborda este estudio es la disponibilidad limitada de información molecular y genómica multipatógeno en aguas servidas de Quito. El trabajo toma como punto de partida la PTAR Quitumbe, caracterizada mediante muestreo pasivo y secuenciación de SARS-CoV-2, y amplía el análisis hacia poliovirus y RSV por su relevancia en vigilancia ambiental, respiratoria y de patógenos con circulación comunitaria \cite{MejiaCalle_2024,Asghar_2014,Hughes_2022,Le_2025}. La detección molecular y la caracterización genómica de las señales virales se interpretan considerando las posibilidades y limitaciones de estas metodologías en matrices ambientales \cite{Child_2023,Parkins_2023}.

Con base en ese problema, la pregunta de investigación es: ¿qué evidencia molecular y genómica de SARS-CoV-2, poliovirus y virus sincitial respiratorio puede identificarse en muestras de aguas servidas recolectadas en la Planta de Tratamiento de Aguas Residuales Quitumbe, en Quito?

El objetivo general es analizar la presencia molecular y la caracterización genómica de SARS-CoV-2, poliovirus y virus sincitial respiratorio en muestras de aguas servidas de la PTAR Quitumbe, como aporte a la vigilancia viral ambiental en Quito.

Los objetivos específicos son:
\begin{enumerate}[label=\arabic*.]
  \item Detectar material genético de SARS-CoV-2, poliovirus y virus sincitial respiratorio en muestras de aguas servidas recolectadas en la PTAR Quitumbe.
  \item Caracterizar, cuando la calidad de los datos lo permita, las señales genómicas virales obtenidas mediante herramientas de secuenciación y análisis bioinformático.
  \item Interpretar los resultados moleculares y genómicos en relación con la vigilancia basada en aguas residuales y el contexto epidemiológico local.
  \item Identificar alcances y limitaciones del análisis multipatógeno en una matriz ambiental compleja como las aguas servidas urbanas.
\end{enumerate}

\subsection{Propósito y estructura del estudio}

El propósito de este estudio es producir evidencia local que permita contrastar el problema de investigación planteado: la limitada disponibilidad de información molecular y genómica multipatógeno en aguas servidas de Quito. Para ello, el trabajo articula detección molecular, caracterización genómica e interpretación bioinformática de SARS-CoV-2, poliovirus y virus sincitial respiratorio en muestras de la PTAR Quitumbe, un punto de muestreo previamente documentado para vigilancia ambiental de SARS-CoV-2 \cite{MejiaCalle_2024,Child_2023}.

El cuerpo del documento se organiza en seis secciones principales. La introducción presenta los antecedentes, la problemática, la justificación, el problema de investigación, la pregunta y los objetivos. El contexto y marco teórico desarrolla la vigilancia epidemiológica, la vigilancia basada en aguas residuales, los virus de interés, las técnicas de detección molecular, la secuenciación, las herramientas bioinformáticas y el contexto epidemiológico del estudio. La metodología y diseño de la investigación describe la justificación metodológica, las herramientas utilizadas, las muestras y los procedimientos de recolección y análisis de datos. El análisis de datos presenta los resultados moleculares y genómicos obtenidos. La discusión interpreta esos resultados frente a la literatura y a la pregunta de investigación. Las conclusiones responden la pregunta, declaran las limitaciones y plantean la importancia y recomendaciones derivadas del estudio. Finalmente, el documento incluye referencias y anexos como soporte académico y técnico.
